% arara: pdflatex: { shell: yes, interaction: nonstopmode }
% arara: pythontex: {verbose: yes, rerun: modified }
% arara: pdflatex: { shell: yes, interaction: nonstopmode }

\documentclass{article}
\usepackage{pythontex}
\usepackage{pgfplots}

\title{Example graph}
\author{Nicky}
\begin{document}
\maketitle

If you want to run the code, but don't want to show it, use the \verb+pycode+ environment.

\begin{figure}[h]
\centering
\begin{pycode}

from matplotlib.pylab import plt
import tikzplotlib

plt.plot([1, 2, 3], [1, 5, 2], label="x")

plt.legend()
print(tikzplotlib.get_tikz_code(axis_height="5cm", axis_width="6cm"))
plt.clf()
\end{pycode}
\caption{It works.}
\end{figure}

If you like to include the code to show how you made the graph, use the \verb+pyblock+ environment, and then use \verb+\printpythontex+ within a figure environment. Like this: 

\begin{pyblock}

from matplotlib.pylab import plt
import tikzplotlib

plt.plot([1, 2, 3], [9, -1, 8], label="x")

plt.legend()
print(tikzplotlib.get_tikz_code(axis_height="5cm", axis_width="6cm"))
\end{pyblock}

\begin{figure}[h]
  \centering
  \printpythontex
\end{figure}

\end{document}
